% Full exact masks, smooth/rough proof and elementary quantitative bounds.
\section{Proof of the rough-prime mask budget}\label{rma:details}
\subsection{The common ideal divisor and the exact rough part}
Here~\eqref{sd:form} uses the complete original good squarefree
annulus and common $W$, without an additional selected-column or
column-ray indicator. The actual row filters are unchanged.
The identity $\mathbf1_{(m,n)=1}=\sum_{d\mid m,\ d\mid n}\mu(d)$
in~\eqref{sd:form} is an identity of ideal divisors. Write $m=da,n=db$.
The residual ideals $a,b$ are squarefree but may overlap; squarefreeness
of the original columns requires $(d,ab)=1$. On $(d,u)=1$ the original
common phases satisfy
\[
 |\nu(d)|^2\chi_d(u)\overline{\chi_d(u)}=1.
\]
On $(d,u)>1$ the product is the original raw zero. Hence the full deleted
set is $\mathcal T=S\cup\operatorname{supp}(abu)$, including residual
primes where $\chi_a\overline{\chi_b}$ might become imprimitive.
With
\[
 F_{a,b}(x)=W(xNa/D)\overline{W(xNb/D)},\qquad
 M=\max(Na,Nb),\quad\alpha=\min(Na,Nb),
\]
the inner signed sum is
\begin{equation}\label{rma:inner}
 \sum_{\substack{d\ \mathrm{SF}\\(d,\mathcal T)=1}}
             \mu(d)A_U(da)A_U(db)F_{a,b}(Nd).
\end{equation}
If its profile is nonzero, then
\begin{equation}\label{rma:support}
 D/\alpha<Nd<2D/M.
\end{equation}
All endpoints have zero profile value. In particular the profile is
identically zero if $\alpha\le M/2$. No independent marginal row family
or changed profile has been substituted.

Factor $d=fe$ uniquely, with $f$ squarefree $U$-smooth and $e$ squarefree
with every prime norm greater than $U$. Whenever $ez$ is squarefree,
\begin{equation}\label{rma:invariance}
 A_U(ez)=A_U(z).
\end{equation}
Indeed neither short divisor in~\eqref{rma:AU} can contain a prime of
$e$; deleting $e$ preserves the two cutoffs and both M\"obius signs in
a bijection of the original assignments. There is no factor $\mu(e)$
in~\eqref{rma:invariance}. This invariance is specific to $A_U$; it does
not hold unchanged for the other coefficient channels. The values
$A_U(fa),A_U(fb)$ are not factorized.

If $M>1$, goodness gives $M\ge7$. From~\eqref{rma:window} and
\eqref{rma:support}, $Nd\le2D/M\le U^2$. Thus $e$ is either $1$ or
one prime ideal, and~\eqref{rma:inner} is exactly
\begin{equation}\label{rma:rough-inner}
 \sum_{\substack{f\ \mathrm{good\ SF},\ U\text{-smooth}\\
                  (f,\mathcal T)=1}}
 \mu(f)A_U(fa)A_U(fb)
 \left(F_{a,b}(Nf)-
       \sum_{\substack{p\notin\mathcal T\\Np>U}}
                       F_{a,b}(NfNp)\right).
\end{equation}
The rough part here belongs to the \emph{common divisor} $d$; the full
columns can have additional rough primes in $a$ or $b$.

If $M=1$, then $a=b=1$. Three rough prime ideals would have norm product
strictly greater than $7U^2\ge2D$: two contribute more than $U^2$, and
the third good prime has norm at least $7$. Two rough primes are
possible, but then a nonunit $f$ would similarly give
$Nf\,Ne>7U^2\ge2D$. Hence that case has $f=1$ only. Its sign is
$\mu(pq)=1$ and its coefficient is $A_U(pq)=A_U(1)=1$.
Summing this case over the unchanged literal rows gives
\eqref{rma:Ctwo}; summing all other cases of
\eqref{rma:rough-inner} over $a,b$ gives the completed squares in
\eqref{rma:Kzero}--\eqref{rma:KT}. This proves~\eqref{rma:exact}.
The unit ideal among residuals and conjugate split primes must both be
retained. The condition $f=1$ is not a termwise assertion that the
original pair has gcd $1$: the M\"obius expansion cancels noncoprime
original pairs across its divisor sectors.

\subsection{An exact disjoint error identity}
We restore only $(p,abuS)=1$ to $p\notin S$. In particular the conditions
$(f,abuS)=1$ and every raw $\chi_a(u)$ remain in the two residual sums.
The full deleted union has the disjoint partition
\begin{equation}\label{rma:union}
 \mathbf1_{p\mid abu}=\mathbf1_{p\mid u}
       +\mathbf1_{p\nmid u}\mathbf1_{p\mid a\ \mathrm{or}\ p\mid b}.
\end{equation}
Therefore $K_S-K_T=E_{\rm row}+E_{\rm res}$, where the errors are
\emph{unnormalized}:
\begin{align}
 E_{\rm row}
 &=\sum_u w_u\sum_{\substack{f\ \mathrm{good\ SF},\ U\text{-smooth}\\
                            (f,u)=1}}\mu(f)
        \sum_{\substack{p\notin S,\ Np>U\\p\mid u}}
                  |T_f(Np;u)|^2,\label{rma:Erow}\\
 E_{\rm res}
 &=\sum_u w_u\sum_{\substack{f\ \mathrm{good\ SF},\ U\text{-smooth}\\
                            (f,u)=1}}\mu(f)
        \sum_{\substack{p\notin S,\ Np>U\\p\nmid u}}
       \bigl(|T_f(Np;u)|^2-|T_f^{(p)}(Np;u)|^2\bigr).
       \label{rma:Eres}
\end{align}
The two errors are signed; their bounds below use absolute values,
and neither is assumed nonnegative.
At $p\mid u$, all residuals divisible by $p$ already have raw character
zero. This is consistent with the full square in~\eqref{rma:Erow}; it
does not make that square zero. These formulas keep the complete
row/residual union without double counting.

Suppose a nonzero omitted residual has $p\mid a$, and write $a=pc$.
Its profile requires $Nf\,(Np)^2Nc<2D$. If $fc$ is nonunit, then
\[
 Nf\,(Np)^2Nc>7U^2\ge2D,
\]
a contradiction. Thus necessarily $f=c=1$ and $a=p$. This is a support
statement about the explicitly auxiliary extension; the formal column
$fpa=p^2$ is not an original squarefree column. With $P=Np$, set
\begin{equation}\label{rma:z}
 z_p(u)=\nu(p)\chi_p(u)W(P^2/D),\qquad T_p(u)=T_1(P;u).
\end{equation}
Here $A_U(p)=1$, since $P>U$. The difference $T_f-T_f^{(p)}$ vanishes
unless $f=1$, when it equals $z_p$. Consequently
\begin{equation}\label{rma:res-exact}
 E_{\rm res}=\sum_u w_u
       \sum_{\substack{p\notin S,\ Np>U\\p\nmid u}}
       \left(2\operatorname{Re}(T_p(u)\overline{z_p(u)})
                         -|z_p(u)|^2\right).
\end{equation}
The subtraction $-|z_p|^2$ is the exact intersection
$p\mid a$ and $p\mid b$; omitting it would count that intersection
twice. Complex profiles and all remaining literal phases are allowed.
If $z_p\ne0$, then
\begin{equation}\label{rma:zsupport}
 \sqrt D<P<\sqrt{2D}.
\end{equation}
The condition $P>U$ is automatic here from $U^2\le D$.

\subsection{Elementary ideal and coefficient estimates}
We give all counting inputs; no prime distribution theorem or Mertens
estimate is needed. The Eisenstein norm is Euclidean. After rounding
both real coordinates of a complex quotient in the basis $1,\omega$,
the error $s+t\omega$ has $|s|,|t|\le1/2$ and norm at most $3/4<1$.
A least-norm nonzero element of any ideal therefore generates it by
Euclidean division. Multiplication by $a+b\omega$ has determinant
$a^2-ab+b^2$, giving the principal ideal norm. Since
$a^2-ab+b^2\ge(a^2+b^2)/2$, counting all possible generators, with
harmless overcounting by units, proves for $X\ge1$ that
\begin{equation}\label{rma:lattice}
 \#\{a:Na\le X\}\le(2\sqrt{2X}+1)^2<16X.
\end{equation}
An annulus $L<Na<2L$ therefore contains at most $32L$ ideals when
$2L>1$. If $2L\le1$ its profile contribution is zero, including the
unit endpoint when $2L=1$. This proves the needed uniform bound even
when $1/2<L<1$; the unit ideal must not be dropped in that range.

There is an exact coefficient identity on the restricted support:
\begin{equation}\label{rma:coefficient}
 A_U(g)=\mu(g)\bigl(2\mathbf1_{Ng\le U}-1\bigr)
 \quad(g\ \mathrm{good\ squarefree},\ Ng\le7U).
\end{equation}
Every proper divisor of a nonunit good $g$ has norm at most $Ng/7\le U$.
If $Ng\le U$, all ternary assignments are allowed and their sum is
$\prod_{p\mid g}(1-1-1)=\mu(g)$. If $Ng>U$, only the assignments
$(x,y)=(g,1),(1,g)$ are omitted; each has sign $\mu(g)$, leaving
$-\mu(g)$. The unit ideal gives $A_U(1)=1$ and is included.
No approximation or multiplicativity of the cutoff coefficient is used.

In an active term of~\eqref{rma:T}, $P=Np>U$ gives
\[
 N(fa)<2D/P<2D/U\le7U.
\]
Therefore $|A_U(fa)|=1$ exactly there. This statement concerns only the
rough-prime support, not the entire original coefficient family.
With $L=D/(FP)$, the ideal count and $|\chi_a(u)|\le1$ now give
\begin{equation}\label{rma:Tbound}
 |T_f(P;u)|\le32B\,\frac{D}{FP}.
\end{equation}
The same statement holds if the sum is empty, and covers $L<1$ by the
preceding endpoint discussion. A dyadic use of~\eqref{rma:lattice} gives
\begin{equation}\label{rma:fsum}
 \sum_f\frac1{(Nf)^2}
 \le\sum_{j\ge0}16\,2^{j+1}2^{-2j}=64,
\end{equation}
where the left side may even include all ideals. The smooth squarefree
restriction only decreases this bound.

\subsection{The normalized row-mask budget}
Apply triangle only to~\eqref{rma:Erow}, use~\eqref{rma:Tbound},
and then~\eqref{rma:fsum}. The raw square contributes
$D^2/P^2$; dividing by $D$ once gives
\begin{equation}\label{rma:row-bound}
 \frac{|E_{\rm row}|}{D}
 \le65536B^2
       \sum_u w_u\sum_{\substack{p\mid u,\ p\notin S\\Np>U}}
                                      \frac{D}{(Np)^2}.
\end{equation}
By~\eqref{rma:window}, $D/(Np)^2\le D/U^2\le7/2$.
Let $\kappa_U(u)$ count the distinct good rough prime ideals dividing $u$.
Each $P=Np$ is at least $7$ and greater than $U$, so $P^2>7U\ge2U$.
The complete ideal norm of the literal row therefore gives
\begin{equation}\label{rma:log}
 \kappa_U(u)\le\frac{2\log(Nu)}{\log(2U)}
          \le\frac{2\log(2H)}{\log(2U)}.
\end{equation}
The denominator is positive also at $U=1$; unit rows have $\kappa_U=0$.
All sixth-factor divisors remain in this bound, and no rational-norm gcd
is substituted for an ideal mask. Using $D/(Np)^2\le7/2$ in~\eqref{rma:row-bound} gives the sharper form
\begin{equation}\label{rma:row-final}
 \frac{|E_{\rm row}|}{D}\le229376B^2
                              \sum_u w_u\kappa_U(u).
\end{equation}
The logarithmic bound~\eqref{rma:log} is available for arbitrary $H$.

\subsection{The residual intersection budget}
Use~\eqref{rma:res-exact}, $|z_p|\le B$, and
\eqref{rma:Tbound} at $f=1$. After the single division by $D$,
\begin{equation}\label{rma:res-bound}
 \frac{|E_{\rm res}|}{D}
 \le B^2R_w\sum_{\sqrt D<Np<\sqrt{2D}}
       \left(\frac{64}{Np}+\frac1D\right).
\end{equation}
Here absolute values bound only the correction; the exact identity
\eqref{rma:res-exact} retains its original phases. Since $D\ge1$ and
$Np\le\sqrt{2D}\le2D$, one has $1/D\le2/(Np)$. Also
\begin{equation}\label{rma:reciprocal}
 \sum_{\sqrt D<Np<\sqrt{2D}}\frac1{Np}\le16\sqrt2
\end{equation}
by~\eqref{rma:lattice}. Thus
\[
 |E_{\rm res}|/D\le1056\sqrt2\,B^2R_w\le2048B^2R_w.
\]
Together with~\eqref{rma:row-final}, this proves~\eqref{rma:delta}.

\subsection{The exact unit-residual correction and the clean main}
In~\eqref{rma:Ctwo}, each contributing pair has $P,Q>U$ and $PQ\le2D$.
Hence both prime norms are less than $2D/U$. There are at most
$32D/U$ such prime ideals by~\eqref{rma:lattice}, so the number of
unordered distinct pairs is at most $512D^2/U^2$. Only inside this
\emph{nonnegative} expression enlarge the masked row count to $R_w$
and the profile to $B^2$. Thus
\[
 0\le\frac{C_2}{D}\le512\frac{D}{U^2}B^2R_w
                         \le1792B^2R_w,
\]
proving~\eqref{rma:Cbound}. Conjugate primes of equal rational norm
remain included; no rational-squarefree restriction or new diagonal
exclusion is imposed. This is a bookkeeping bound for the exact
correction, not an oscillatory saving or a lower bound for the full
signed slice. Other residual sectors can cancel it.

Now~\eqref{rma:exact} gives~\eqref{rma:clean} with
$E=\Delta+C_2/D$, and the explicit bounds give~\eqref{rma:error}.
Suppose $H\le D^{3/4}$. If two distinct good rough primes divided a row,
then $H\ge Nu\ge(Np)(Nq)\ge49$, but also
\[
 H\ge Nu\ge(Np)(Nq)>U^2\ge2D/7.
\]
Combining with $H\le D^{3/4}$ gives $D^{1/4}<7/2$ and hence
$H<(7/2)^3=343/8<49$, a contradiction. Thus $\kappa_U(u)\le1$
at every scale in this capacity range. Substitution in~\eqref{rma:error}
gives $|E|\le233216B^2R_w\le2^{18}B^2R_w$, proving
\eqref{rma:critical-error}. This is an exact endpoint argument, not an
asymptotic approximation.

For a two-pivot block, counting each of its two dyadic ideal variables
by~\eqref{rma:lattice} gives $R_w\ll A_1A_2$ when the original fixed unit
classes and multiplicities are bounded. Internal coprimality and capacity
reduce that count; grouping literal fibers preserves total multiplicity.
Under $A_1A_2\le H\le D^{3/4}$ the correction is therefore $O(B^2H)$.
The exact local coefficient identity and these correction bounds do not
estimate the remaining signed main $(K_0-K_S)/D$. Its original
$A_U(fa)$, coherent rows, common fixed profile and M\"obius mean remain
the analytic obligation; no new global exponent follows.


\subsection{Finite control scope}
The tables in the shared core retain complete original squarefree ideal
columns and raw character zeros. Their residual-prime extension is a
separate auxiliary sum, including the non-squarefree $p^2$ terms needed
for~\eqref{rma:res-exact}; those are not original columns. The exact
intersection subtraction and unit residual states with $1/2<L<1$ are
checked separately. The common profile has fixed $\delta=1/4096$ and
large fixed derivative constants. The capacities and row blocks differ
between cases, and the $H=384$ case is outside the critical capacity
range. These finite equalities neither determine the analytic constants
nor imply an asymptotic saving.

\subsection{A coherent prime-substitution budget, and its remaining cost}
This subsection requires $U\ge7$. It does not extend the prime-density
argument to $U=1$, although the coefficient and mask lemmas above hold
for all $U\ge1$. For real $t\ge U$, put
\[
 G_f(t)=\sum_{\substack{u\in\mathcal R\\(f,u)=1}}
                      w_u|T_f(t;u)|^2,
 \qquad M_1(W)=\max\{B,\|xW'(x)\|_\infty\}.
\]
The active residuals satisfy $Na\le2D/(FU)$, so their union is finite;
$G_f$ is smooth and vanishes for $t\ge2D/F$. Since the rough prime in
$K_S$ is unmasked, it enters only this radial common profile. Let
\[
 E_S(t)=\pi_K^S(t)-\operatorname{Li}(t),\qquad
 \delta_S(U)=\sup_{t\ge U}\frac{|E_S(t)|}{t},
 \qquad
 \mathcal J=\sum_f\mu(f)\int_U^\infty\frac{G_f(t)}{\log t}\,dt,
\]
where $\pi_K^S$ counts prime \emph{ideals} outside the fixed finite $S$ with
multiplicity, including inert prime norms and distinct conjugate ideals. Stieltjes partial
summation gives the exact coherent identity
\begin{equation}\label{rma:Stieltjes}
 \begin{aligned}
 K_S&=\mathcal J+\mathcal E_{\rm PIT},\\
 \mathcal E_{\rm PIT}
 &=\sum_f\mu(f)\left(-G_f(U)E_S(U)
          -\int_U^\infty E_S(t)G_f'(t)\,dt\right).
 \end{aligned}
\end{equation}
The lower term uses $\pi_K^S(U)$, so the condition $Np>U$ is exact even
when a prime norm equals $U$. A conjugate pair contributes a jump of
two. This equality precedes any absolute sum over $f$.

Let $\lambda(L)$ be a TRUE arbitrary-coefficient squared norm bound for
the specified positive row operator on good squarefree ideals of norm
in $[L,2L]$, including the actual weights. For $1/2\le L<1$, set
$\lambda_*(L)=R_w$ by the direct unit-column bound; for $L\ge1$, use
$\lambda_*(L)=\lambda(L)$. No columns are active for $L<1/2$.
With $L=D/(Ft)$, differentiate only the common $W$: multiplication by
$t\,d/dt$ replaces $W(x)$ by $xW'(x)$. The exact restricted identity
\eqref{rma:coefficient}, ideal counting and weighted row Cauchy give
\[
 G_f(t)+|tG_f'(t)|\ll M_1(W)^2\lambda_*(L)L.
\]
Any enlargement of a common $f$-row mask here occurs only inside these
positive norm estimates. Substitution in~\eqref{rma:Stieltjes}, with
$L_f=D/(FU)$, proves
\begin{equation}\label{rma:PIT-variation}
 \frac{|\mathcal E_{\rm PIT}|}{D}
 \ll M_1(W)^2\delta_S(U)
 \sum_{\substack{f\ \mathrm{good\ SF},\ U\text{-smooth}\\F<2D/U}}
 \frac1F\left(\lambda_*(L_f)
       +\int_{1/2}^{L_f}\lambda_*(L)\,\frac{dL}{L}\right).
\end{equation}
The original M\"obius sign remains in $K_S$ and $\mathcal J$ themselves;
the absolute $f$ sum appears only in this error majorant.

Specialize now ONLY to the established positive input for the actual $A_6=1$
unweighted two-pivot block, or bounded natural literal weights for which
the same input has been established. Put
\[
 R=A_1A_2,\qquad M_{\rm row}=\max(A_1,A_2),\qquad A_1A_2^2\le H.
\]
At $L\ge1$ that input is
\begin{equation}\label{rma:positive-input}
 \lambda(L)\ll_\epsilon(HD)^\epsilon
 \min\left\{R+\frac{LR}{M_{\rm row}}
       +\frac{L^{2/3}R}{M_{\rm row}^{1/3}},\ L+R^2\right\}.
\end{equation}
It is not asserted for arbitrary nonnegative weights, enlarged primitive
conductor families or arbitrary sixth-factor enlargements. Using its
first branch in~\eqref{rma:PIT-variation} gives
\begin{align}
 \frac{|\mathcal E_{\rm PIT}|}{D}
 &\ll_\epsilon(HD)^\epsilon M_1(W)^2\delta_S(U)\mathcal B_{\rm op},
                  \label{rma:PIT-budget}\\
 \mathcal B_{\rm op}
 &=R\log^2(2D)+\frac{D}{U}\frac{R}{M_{\rm row}}
       +\left(\frac{D}{U}\right)^{2/3}
                              \frac{R}{M_{\rm row}^{1/3}}.
                  \label{rma:Bop}
\end{align}
The logarithmic term uses $\sum_{Nf\le X}(Nf)^{-1}\ll\log(2X)$;
the other terms use convergent ideal sums of powers $2$ and $5/3$.
The same positive bounds give
$|\mathcal J|/D\ll_\epsilon(HD)^\epsilon B^2
\mathcal B_{\rm op}/\log U$. These are inherited majorants, not a new
signed gain or an estimate of $K_0-\mathcal J$.

Capacity implies
\[
 \mathcal B_{\rm op}\ll H\log^2(2D)
       +(D/U)H^{1/3}+(D/U)^{2/3}H^{2/3}.
\]
At $H=D^{3/4}$ and $U\asymp\sqrt D$, these terms have powers
$3/4,3/4,5/6$. The last power occurs in this majorant already on pure
scales $A_1\asymp H,A_2=1$ with fixed capacity slack. It is not an
actual-error lower bound.

For the fixed field, Thorner--Zaman,
\href{https://doi.org/10.2140/ant.2019.13.1039}{\emph{Algebra \& Number Theory}
13 (2019), 1039--1068, Theorem~1.1}, applied to the trivial extension
$K/K$ with absolute degree $[K:\mathbb Q]=2$ and absolute discriminant $3$, supplies for sufficiently
large $U$ a bound
\[
 \delta_S(U)\ll_{K,S}\exp(-c_K\sqrt{\log U}).
\]
The possible exceptional real zero is fixed: retain its secondary
$\operatorname{Li}(x^\beta)$ term, or absorb it only into fixed-field
constants and the threshold. Finite $S$ deletions are included. No
prime theorem uniform in the growing deleted set or row conductor is
used. The DOI, publisher URL and pinned reading-cache hashes are in the
\href{https://github.com/amelie-iska/GRH-research/blob/main/outputs/checks/2026-11-05-rough-mask/main-attempt/PUBLIC-CITATION-v1.json}{public citation record}.

A fixed-field subpower prime error does not make the available
$D^{5/6+\epsilon}\delta_S(U)$ majorant certify the target
$D^{3/4+\epsilon}$ for every $\epsilon>0$. Actual signed errors might
cancel. Even an exact prime density would leave the coherent signed
$K_0-\mathcal J$ unestimated. The elementary correction theorem uses
only $B$, whereas the prime-substitution error requires $M_1(W)^2$.
The complete original remainder's $7/6$ envelope remains unchanged.
